\label{app:sql-walkthrough}
This appendix follows one small dataset from insertion to alignment and
index maintenance. Read the left column from top to bottom. SQL statements
are executable, while the differently coloured plan blocks are output from
\texttt{EXPLAIN ANALYZE}. The right column explains each block.

The accompanying \path{appendices/02_sql_walkthrough/run.sql} executes the
same SQL files used to typeset this listing. Run it with
\texttt{psql -X -d bioseq\_walkthrough -f run.sql} in a new, empty local
database created with \texttt{createdb bioseq\_walkthrough}.
The BioSeqDB library must already be installed on the server.
Use a database administrator account for this exercise, since it installs
the extension and writes a small FASTA fixture to
\path{/tmp/bioseq-walkthrough.fasta} on the server.
Each statement commits separately under the default \texttt{psql} behaviour.
Do not wrap the whole script in one transaction.

No alignment GUC is changed in this walkthrough. It uses the defaults
listed in Appendix~\ref{app:configuration-reference}, with explicit
per-call options where they are being demonstrated. Names and dates are
fictional. The DNA is generated deterministically to keep long strings out
of the listing.

\warn{Alpha-version limitations}{%
\label{warn:walkthrough-alpha}
Please note that the current software is an alpha version and still has
numerous issues. Its memory checks can be too restrictive. They keep at
least 2~GiB available before admitting more work and also limit the
process's memory use. Leave several GiB available when trying these examples,
especially when building grouped MPF indexes.

One error-reporting bug can turn a failed query-minimizer allocation into
\texttt{transient query execution stats disagree with its generated backend}.
Free memory and retry before assuming that an index is missing.
Building a stored target index does not remove the need to generate query
minimizers. The message does not establish that the query rejected that
index. Warning~\ref{warn:index-query-minimizers} explains this distinction.

Keep \path{bioseq.index.insert_maintenance.enabled} off, as it is by default.
Ordinary per-row minimizer maintenance is unsupported with the default
grouped format. Build it explicitly after inserting rows, as shown below.
Turning maintenance off does not disable memory checks or repair an
allocation failure.}

A short exact substring need not produce an alignment. Minimizers must
survive seeding and chaining before a hit is reported.
For the fixture below, 30- and 40-base prefixes produce no hits under the
default settings, whereas a 50-base prefix does. This is an example, not a
universal length threshold. The walkthrough uses 256- and 512-base probes.
The seeding and chaining requirements are described in
Sections~\ref{sec:configured-cpu-minimizers} and~\ref{sec:chaining-bioseqdb}.

The displayed plans are complete except for the explicitly marked vector
aggregate excerpt. Times, cache counters, worker counts and CPU instruction
names will vary between runs. They are illustrations of the fields, not
performance benchmarks.
